\documentclass[11pt, a4paper]{article}

\usepackage[T1]{fontenc}
\usepackage[utf8]{inputenc}
\usepackage{tgpagella}
\usepackage[spanish, es-noshorthands]{babel}
\usepackage{amsmath, amssymb, bm}
\usepackage{booktabs}
\usepackage{tabularx}
\usepackage[table, xcdraw]{xcolor}
\usepackage[most]{tcolorbox}
\usepackage[margin=2.5cm]{geometry}
\usepackage{fancyhdr}

\pagestyle{fancy}
\fancyhf{}
\setlength{\headheight}{16pt}
\lhead{\textbf{UCB Sede Tarija} | Ing. Mecatrónica}
\chead{}
\rhead{IMT-121 Circuitos Electrónicos I | Solucionario Docente S06-3}
\lfoot{Prof: M.Sc. Ing. Bernardo Quiroga Turdera}
\rfoot{Página \thepage}
\renewcommand{\headrulewidth}{0.4pt}

\definecolor{ucbblue}{RGB}{0, 51, 102}

\begin{document}

\begin{tcolorbox}[colback=ucbblue!5!white, colframe=ucbblue, title=\textbf{Clave de Evaluación Docente: Thévenin, Norton y Superposición}]
\end{tcolorbox}

\section*{Solución al Ejemplo Guiado (Artefacto 2)}
\textbf{Paso 1: Aislar el Puerto y Hallar $V_{th}$}\\
$V_{th}$ es el voltaje en el nodo $a$ con el puerto abierto. Usando KCL o superposición:
\begin{itemize}
    \item $V_a^{(V_s)} = 12 \cdot \frac{8}{12} = 8\,\text{V}$. \quad $V_a^{(I_s)} = 1\,\text{mA} \cdot (4\,\text{k} \parallel 8\,\text{k}) = 2.667\,\text{V}$.
    \item $\mathbf{V_{th} = 10.667\,\text{V}}$.
\end{itemize}

\noindent\textbf{Paso 2: Desactivar Fuentes y Hallar $R_{th}$}\\
Fuentes de tensión $\rightarrow$ Cortocircuito. Fuentes de corriente $\rightarrow$ Circuito abierto.
Mirando desde $a-b$, los resistores de $4\,\text{k}\Omega$ y $8\,\text{k}\Omega$ quedan en paralelo (ambos están conectados entre el nodo $a$ y tierra).
$\mathbf{R_{th} = 4\,\text{k} \parallel 8\,\text{k} = 2.667\,\text{k}\Omega}$.

\noindent\textbf{Paso 3: Norton y Verificación Final}\\
$R_N = R_{th} = \mathbf{2.667\,\text{k}\Omega}$. \quad $I_N = \frac{10.667\,\text{V}}{2.667\,\text{k}\Omega} = \mathbf{4\,\text{mA}}$. \\
Verificación con $R_L = 8\,\text{k}\Omega$ en la serie de Thévenin:
$I_L = \frac{10.667}{2.667\,\text{k} + 8\,\text{k}} = \mathbf{1\,\text{mA}}$. $V_L = 1\,\text{mA} \cdot 8\,\text{k}\Omega = \mathbf{8\,\text{V}}$.

\section*{Solución al Ejercicio 2 (Artefacto 3)}
\begin{enumerate}
    \item \textbf{$V_{th}$:} Con puerto abierto, por superposición: 
    $V_a^{(1)} = 9 \cdot \frac{6}{9} = 6\,\text{V}$. \quad $V_a^{(2)} = 0.5\,\text{mA} \cdot (3\,\text{k} \parallel 6\,\text{k}) = 0.5\,\text{mA} \cdot 2\,\text{k}\Omega = 1\,\text{V}$.
    $\mathbf{V_{th} = 6 + 1 = 7\,\text{V}}$. \textit{(Error típico: calcular esto con la carga conectada).}
    \item \textbf{$R_{th}$:} Desactivando fuentes: $V_s$ es un corto, $I_s$ es un abierto. Mirando desde $a-b$, se ven $3\,\text{k}\Omega$ y $6\,\text{k}\Omega$ en paralelo. $\mathbf{R_{th} = 2\,\text{k}\Omega}$. \textit{(Error típico: deducir la resistencia sin redibujar mentalmente el corto de la fuente).}
    \item \textbf{$I_N$:} $I_N = \frac{7\,\text{V}}{2\,\text{k}\Omega} = \mathbf{3.5\,\text{mA}}$. \textit{(Error típico: medir $I$ en una rama interna en vez de hacer el corto de puerto).}
    \item \textbf{Verificación:} Usando Thévenin ($7\,\text{V}$ en serie con $2\,\text{k}\Omega$), si conectamos $R_L = 6\,\text{k}\Omega$: \\
    $I_L = \frac{7}{2\,\text{k} + 6\,\text{k}} = 0.875\,\text{mA}$. $\mathbf{V_L = 0.875\,\text{mA} \cdot 6\,\text{k}\Omega = 5.25\,\text{V}}$.
\end{enumerate}

\section*{Solucionario Guía de Práctica S1-S3 (Artefacto 4)}
\begin{itemize}
    \item \textbf{S1:} $V_o^{(1)} = 10(3/5) = 6\,\text{V}$. $V_o^{(2)} = -1\,\text{mA}(1.2\,\text{k}) = -1.2\,\text{V}$. $\mathbf{V_o = 4.8\,\text{V}}$.
    \item \textbf{S2:} $V_o^{(1)} = 12(4/8) = 6\,\text{V}$. $V_o^{(2)} = -2\,\text{mA}(2\,\text{k}) = -4\,\text{V}$. $\mathbf{V_o = 2\,\text{V}}$.
    \item \textbf{S3:} $V_o^{(1)} = 15(10/15) = 10\,\text{V}$. $V_o^{(2)} = +0.5\,\text{mA}(3.333\,\text{k}) = 1.667\,\text{V}$. $\mathbf{V_o = 11.667\,\text{V}}$.
\end{itemize}

\end{document}
